\documentclass[addpoints,12pt]{exam}

\usepackage{amsmath, amsfonts, amssymb}

% --- Header and Footer Setup ---
\pagestyle{headandfoot}
\firstpageheader{Comp Math}{First Homework}{August 31, 2026}
\runningheader{Comp Math}{First Homework, Page \thepage\ of \numpages}{}
\firstpagefooter{}{}{}
\runningfooter{}{}{Page \thepage}

\begin{document}

\begin{center}
    \textbf{\Large Homework \#1} \\
    \vspace{0.1in}
    \textbf{\large Introduction to Proofs}
\end{center}

\vspace{0.3in}

\begin{questions}

% ------------------------------------------------------------
% Question 1
% ------------------------------------------------------------

\question Factor $2^{15}-1=32{,}767$ into a product of two smaller positive integers.

\begin{parts}

    \part
    Earlier in the book, there is a theorem that essentially states
    \[
        xy
        =
        (2^b-1)
        \left(
            1+2^b+2^{2b}+\cdots+2^{(a-1)b}
        \right),
    \]
    where the number can be written as $2^n-1$ and $a$ and $b$ are
    less than $n$.

    I selected $a=5$ and $b=3$, and rewrote $32{,}767$ as
    \[
        (2^3-1)
        \left(
            \sum_{n=0}^{a-1}2^{bn}
        \right).
    \]

    This gives the two positive factors
    \[
        (7)(4681).
    \]

    \part
    Find an integer $x$ such that
    \[
        1<x<2^{32{,}767}-1
    \]
    and $2^{32{,}767}-1$ is divisible by $x$.

    Using the same logic, we can find that one of the factors of this
    number is
    \[
        2^7-1=127.
    \]

\end{parts}

% ------------------------------------------------------------
% Question 2
% ------------------------------------------------------------

\question
Make some conjectures about the values of $n$ for which $3^n-1$ is
prime, or the values of $n$ for which $3^n-2^n$ is prime.

\begin{parts}

    \part
    \[
    \begin{array}{c|r|c|r|c}
    \boldsymbol{n}
        & \boldsymbol{3^n-1}
        & \text{Prime?}
        & \boldsymbol{3^n-2^n}
        & \text{Prime?}
        \\ \hline
    1  & 2      & \checkmark & 1      & \times \\
    2  & 8      & \times     & 5      & \checkmark \\
    3  & 26     & \times     & 19     & \checkmark \\
    4  & 80     & \times     & 65     & \times \\
    5  & 242    & \times     & 211    & \checkmark \\
    6  & 728    & \times     & 665    & \times \\
    7  & 2186   & \times     & 2059   & \checkmark \\
    8  & 6560   & \times     & 6305   & \times \\
    9  & 19682  & \times     & 17707  & \checkmark \\
    10 & 59048  & \times     & 58025  & \times
    \end{array}
    \]

    For all $n>1 , 3^n - 1$ is not a prime  
\end{parts}

\question 
The proof of Theorem 3 gives a method for finding a prime number different from any in a given list of prime numbers
    \begin{parts}
    \part 
    $(2*3*5*7) + 1 = 211$

    \part 
    $(2*3*5*7*11) + 1 = 2311$
    \end{parts}

\question
Find five consecutive integers that are not prime.

    24,25,26,27,28

\question Use the table in Figure I.1 and the discussion on p. 5 to find two more
perfect numbers
    \begin{parts}
       \part 31 and 127 

    \end{parts}

\question The sequence 3, 5, 7 is a list of three prime numbers such that each
pair of adjacent numbers in the list differ by two. Are there any more
such “triplet primes”?
    \begin{parts}
        \part If a number is divisible by 3 then $mod(n,3) = 0$, however, mod is a cylic function and will either be 0,1,2. with leave only 2 other number that can be factors, breaking the sequence of 3. the only sequence that satisfies this is the one given. 
    \end{parts}
\question A pair of distinct positive integers (m, n) is called amicable if the sum
of all positive integers smaller than n that divide n is m, and the sum
of all positive integers smaller than m that divide m is n. Show that
(220, 284) is amicable.
    \begin{parts}
        \part
        220's factors that are less than 220 are 1,2,4,5,10,11,20,22,44,55,110. The factors add up to 284, so it is amicable.
    \end{parts}

\end{questions}
\end{document}